\section{Conic classification before naming the curve}

Consider the general real conic
\begin{equation}
 Ax^2+Bxy+Cy^2+Dx+Ey+F=0.\label{eq:general-conic}
\end{equation}
Its quadratic matrix is
\begin{equation}
 \matr Q=\begin{bmatrix}A&B/2\\B/2&C\end{bmatrix},\qquad
 \det\matr Q=AC-\frac{B^2}{4}.
\end{equation}
Consequently
\begin{equation}
 B^2-4AC=-4\det\matr Q.\label{eq:disc-det}
\end{equation}
This identity is a bridge between classical analytic geometry and linear
algebra.  The left side is the coefficient test traditionally learned for
conics; the right side asks whether the quadratic map preserves area with one
orientation, reverses it, or collapses a direction.  Since
$\det\matr Q$ is also the product of the two eigenvalues:
\begin{itemize}
 \item $B^2-4AC<0$ means equal eigenvalue signs.  A positive or negative
       definite quadratic part is ellipse-like, but the translated level may
       still be a real ellipse, a point, or empty.
 \item $B^2-4AC=0$ means at least one zero eigenvalue.  Linear terms then decide
       between a parabola, parallel lines, one line, or the empty set.
 \item $B^2-4AC>0$ means opposite eigenvalue signs.  The quadratic contribution
       rises in one principal direction and falls in the other, producing a
       hyperbola or a degenerate pair of intersecting lines.
\end{itemize}
Thus the discriminant classifies the quadratic part; reality and degeneracy
also require the lower-order terms.

A third, distinct determinant belongs to the homogeneous conic matrix
\begin{equation}
 \matr C=\begin{bmatrix}
 A&B/2&D/2\\ B/2&C&E/2\\ D/2&E/2&F
 \end{bmatrix},\qquad
 [x\ y\ 1]\matr C[x\ y\ 1]\trans=0.
\end{equation}
Its determinant detects degeneracy of the \emph{whole} conic, not definiteness
of the $2\times2$ quadratic block.  For a centered ellipse
$\vect z\trans\matr Q\vect z-k=0$,
\begin{equation}
 \det\matr C=-k\det\matr Q<0,
\end{equation}
while simultaneously $\det\matr Q>0$ and $B^2-4AC<0$.  The signs differ
because the three expressions answer three different questions; there is no
contradiction.

For a nonsingular symmetric matrix, complete the square:
\begin{align}
q(\vect x)&=\vect x\trans\matr Q\vect x+2\vect d\trans\vect x+f,\\
\vect x_0&=-\matr Q\inv\vect d,\label{eq:center-general}\\
q(\vect x)&=(\vect x-\vect x_0)\trans\matr Q(\vect x-\vect x_0)
 +f-\vect d\trans\matr Q\inv\vect d.\label{eq:matrix-square}
\end{align}
The cancellation follows from $\matr Q\vect x_0=-\vect d$ and symmetry.  For
the voltage inequality define
\begin{equation}
 k\defeq\vect b\trans\Hpsm\inv\vect b-c.
\end{equation}
Then its boundary is
\begin{equation}
 (\current-\ic)\trans\Hpsm(\current-\ic)=k,
 \qquad \ic=-\Hpsm\inv\vect b.\label{eq:centered-voltage}
\end{equation}
Because $\vect r$ lies in the range of invertible $\Apsm$, the current
$\ic=-\Apsm\inv\vect r$ maps to zero voltage.  Hence directly $k=\Umax^2>0$.
We have now established all conditions for a real, nondegenerate ellipse:
$\Hpsm\succ0$ and $k>0$.

Explicit inversion yields
\begin{equation}
 \boxed{\ic=\begin{bmatrix}
 -\dfrac{\psif\omegae^2\Lq}{\Rs^2+\omegae^2\Ld\Lq}\\[.7em]
 -\dfrac{\psif\Rs\omegae}{\Rs^2+\omegae^2\Ld\Lq}
 \end{bmatrix}.}\label{eq:ellipse-center}
\end{equation}

\begin{figure}[tb]
\centering
\begin{tikzpicture}[scale=.92]
 \draw[->] (-3.8,0)--(3.4,0) node[right]{$i_d$};
 \draw[->] (0,-2.1)--(0,2.4) node[above]{$i_q$};
 \draw[densely dashed,gray,thick] (0,0) ellipse (1.8 and 1.05);
 \draw[constraint curve,rotate around={12:(-1.35,.35)}]
   (-1.35,.35) ellipse (1.8 and 1.05);
 \draw[coordinate vector] (0,0)--(-1.35,.35) node[midway,above]{$\ic$};
 \node[operating point] at (0,0) {};
 \node[psm center] at (-1.35,.35) {};
 \node[annotation] at (1.9,-1.5){translation: center\\rotation: orientation};
\end{tikzpicture}
\caption{Completing the square translates the center.  Translation alone does
not change the conic type or semi-axis lengths.}
\label{fig:translation}
\end{figure}


\paragraph{Why complete the square?}
The linear term prevents the quadratic part from being read as a shape around
its natural center.  Completing the square is therefore not cosmetic algebra:
it asks where the quadratic landscape has its stationary point and translates
that point to the origin.  Translation exposes the center and the remaining
energy level, but it deliberately hides the absolute current origin.  Directly
solving $\Apsm\current+\vect r=\vect0$ is an alternative route to the same
center and makes the zero-voltage interpretation more visible.
